%! Inverse Functions — Unit 1, Lesson 1.6 — Guided Notes (Answer Key)
\documentclass[10pt]{article}
\usepackage{precalculus-article}
\usepackage{precalculus-key}   % pulls in -boxes; do NOT also load -boxes
\newcommand{\TallMath}[1]{$\displaystyle #1\rule[-1.4em]{0pt}{3.2em}$}

\begin{document}

\pageheader{Unit 1, Lesson 1.6}{Guided Notes --- Answer Key}

\vspace{0.12in}

\begin{objectivebox}
By the end of this lesson, I will be able to\ldots
\begin{enumerate}[itemsep=2pt]
  \item Reverse the input--output pairs of a function to build $f^{-1}$, and
        explain why the \ans{domain} and \ans{range} trade places.
  \item Find $f^{-1}(x)$ by writing $y = f(x)$, \ans{swapping} $x$ and $y$, and
        \ans{solving} for $y$.
  \item Verify a pair of inverses using \ans{composition}, and restrict a domain
        when a function is not \ans{one-to-one}.
\end{enumerate}
\end{objectivebox}

\vspace{0.1in}

\boxguard
\begin{vocabbox}
Fill in each term as we define it together.
\par\vspace{2pt}
\noindent\textbf{\textcolor{plum}{Inverse function:}}\\[1pt]
\ansline{The function that reverses $f$: if $f(a) = b$, then $f^{-1}(b) = a$.}\\[3pt]
\noindent\textbf{\textcolor{plum}{One-to-one:}}\\[1pt]
\ansline{No two different inputs give the same output --- outputs never repeat.}\\[3pt]
\noindent\textbf{\textcolor{plum}{Domain restriction:}}\\[1pt]
\ansline{Cutting the domain down to a piece where the function is one-to-one.}\\[3pt]
\noindent\textbf{\textcolor{plum}{Identity function:}}\\[1pt]
\ansline{The function that returns its own input, $I(x) = x$.}\\[3pt]
\noindent\textbf{\textcolor{plum}{Reflection over $y = x$:}}\\[1pt]
\ansline{$f^{-1}$'s graph is the mirror image of $f$'s across the line $y = x$.}
\end{vocabbox}

\vspace{0.1in}

\boxguard
\begin{hookbox}
The weather app says $68^\circ$F. Your friend in Paris wants that in Celsius.
The only rule on the board is $F(x) = 1.8x + 32$ --- and it converts Celsius
\textit{into} Fahrenheit, which is the wrong direction.

\smallskip
What do you do to $68$ to get back to Celsius, and in what \textit{order}?

\ansline{Undo the rule backwards: subtract 32 first, then divide by 1.8.}
\ansline{$(68 - 32) \div 1.8 = 20$, so $20^\circ$C. That undoing is itself a function.}
\end{hookbox}

\vspace{0.1in}

\boxguard[26]
\begin{notesbox}{1. ``Undoing'' Is a Reversed Mapping}
The rule $F(x) = 1.8x + 32$ takes a Celsius temperature and returns a
Fahrenheit one. Its \textbf{inverse}, written $F^{-1}$, runs the same pairs
backwards.

\medskip
\begin{center}
\renewcommand{\arraystretch}{1.6}
\begin{tabularx}{0.9\linewidth}{@{} >{\bfseries}p{0.32\linewidth} X @{}}
\hline
$F$ sends \ldots\ to \ldots & $F^{-1}$ sends it back \\
\hline
$F(20) = 68$ & $F^{-1}(68) = $ \ans{20} \\
\hline
$F(0) = 32$  & $F^{-1}(32) = $ \ans{0} \\
\hline
$F(a) = b$   & $F^{-1}($ \ans{$b$} $) = $ \ans{$a$} \\
\hline
\end{tabularx}
\end{center}

\medskip
\textbf{The definition, in one line:} if $f(a) = b$, then $f^{-1}(b) = a$.

\medskip
\textbf{Warning --- the biggest mistake in this lesson.} The $-1$ in $f^{-1}$
is a \textbf{name}, not an exponent:
\begin{center}
$f^{-1}(x) \ne \dfrac{1}{f(x)}$
\end{center}

See it in numbers. Since $F(10) = 1.8(10) + 32 = 50$, we get
$F^{-1}(50) = $ \ans{10}. \\
But $F(50) = 1.8(50) + 32 = $ \ans{122}, so
$\dfrac{1}{F(50)} = $ \ans{$1/122$}.

\smallskip
Those two answers are \ans{nowhere near each other}.
\end{notesbox}

\vspace{0.1in}

\boxguard[26]
\begin{notesbox}{2. Inverses from Tables --- and the Domain/Range Swap}
\textbf{Example 1.} Complete the Fahrenheit row for $F(x) = 1.8x + 32$.

\smallskip
\renewcommand{\arraystretch}{1.4}
\begin{center}
\begin{tabular}{|c|c|c|c|c|c|}
\hline
$x$ (degrees C) & $-10$ & 0 & 10 & 20 & 100 \\
\hline
$F(x)$ (degrees F) & \ans{14} & \ans{32} & \ans{50}
  & \ans{68} & \ans{212} \\
\hline
\end{tabular}
\end{center}

\smallskip
Now build $F^{-1}$ by \textbf{trading the two rows}. Nothing is computed here
--- the same pairs are read backwards.

\smallskip
\renewcommand{\arraystretch}{1.4}
\begin{center}
\begin{tabular}{|c|c|c|c|c|c|}
\hline
$x$ (degrees F) & 14 & 32 & 50 & 68 & 212 \\
\hline
$F^{-1}(x)$ (degrees C) & \ans{$-10$} & \ans{0} & \ans{10}
  & \ans{20} & \ans{100} \\
\hline
\end{tabular}
\end{center}

\medskip
\begin{itemize}[itemsep=4pt]
  \item Domain of $F^{-1}$ $=$ the \ans{range} of $F$. \quad
        Range of $F^{-1}$ $=$ the \ans{domain} of $F$.
  \item \textbf{The catch.} Reversing works only when no two inputs share an
        output. Take $g = \{(1,5), (2,5), (3,8)\}$. Reversed, the input 5 would
        have to produce \ans{two} different outputs, so the reversal is
        \ans{not a function}.
  \item A function whose outputs never repeat is called \ans{one-to-one}, and
        only such a function is \ans{invertible}.
\end{itemize}
\end{notesbox}

\vspace{0.1in}

\boxguard[26]
\begin{notesbox}{3. The Analytic Method: Swap and Solve}
\begin{enumerate}[itemsep=3pt]
  \item Write $y = f(x)$.
  \item \ans{Swap} $x$ and $y$. \quad \textit{(This is the inversion. The
        rest is algebra.)}
  \item \ans{Solve} for $y$.
  \item Rename the result $f^{-1}(x)$.
\end{enumerate}

\medskip
\textbf{Example 2.} Find $f^{-1}(x)$ for $f(x) = 4x - 1$.

\begin{work}
              y &= 4x - 1 \\
              x &= 4y - 1 \\
          x + 1 &= 4y \\
  \frac{x+1}{4} &= y
\end{work}

$f^{-1}(x) = $ \ans{$(x+1)/4$}

\medskip
\textbf{Example 3.} Find $F^{-1}(x)$ for the temperature rule
$F(x) = 1.8x + 32$.

\begin{work}
                  y &= 1.8x + 32 \\
                  x &= 1.8y + 32 \\
             x - 32 &= 1.8y \\
  \frac{x - 32}{1.8} &= y
\end{work}

$F^{-1}(x) = $ \ans{$(x-32)/1.8$}

\smallskip
Check it against the table: $F^{-1}(68) = \dfrac{68 - 32}{1.8} = $
\ans{20}, which matches the \ans{table} from Part 2.
\end{notesbox}

\vspace{0.1in}

\boxguard[26]
\begin{notesbox}{4. Verification by Composition}
Lesson 1.5 built the tool: $(f \circ g)(x) = f(g(x))$ means run $g$ first, then
feed its output into $f$. Here that tool becomes a \textbf{test}.

\medskip
\begin{center}
$f$ and $f^{-1}$ are inverses exactly when \quad
$f\big(f^{-1}(x)\big) = f^{-1}\big(f(x)\big) = x$.
\end{center}

The function that returns its own input, $I(x) = x$, is the
\ans{identity} function. Composing a function with its inverse --- in
\textbf{either} order --- produces it.

\medskip
\textbf{Example 4.} Verify that $f(x) = 4x - 1$ and $f^{-1}(x) = \dfrac{x+1}{4}$
are inverses.

\begin{work}
  f\big(f^{-1}(x)\big) &= 4\left(\frac{x+1}{4}\right) - 1 \\
                       &= (x + 1) - 1 \\
                       &= x \\
  f^{-1}\big(f(x)\big) &= \frac{(4x - 1) + 1}{4} \\
                       &= \frac{4x}{4} \\
                       &= x
\end{work}

Both compositions collapsed to \ans{$x$}, so $f$ and $f^{-1}$ really are
\ans{inverses}.

\medskip
\textbf{You already did this.} In the warm-up you found $f(g(7)) = 7$ and
$g(f(7)) = 7$ with these exact two functions. The algebra above does the same
check for \ans{every input} at once, not just for 7.

\smallskip
\textbf{Why both directions, and not just one?}

\ansline{Both must collapse to $x$; one direction alone is not enough.}
\end{notesbox}

\vspace{0.1in}

\boxguard[30]
\begin{notesbox}{5. When There Is No Inverse --- and the Mirror Line $y = x$}
\textbf{Example 5.} $f(x) = x^2$ on all real numbers.

\smallskip
\renewcommand{\arraystretch}{1.4}
\begin{center}
\begin{tabular}{|c|c|c|c|c|c|}
\hline
$x$ & $-2$ & $-1$ & 0 & 1 & 2 \\
\hline
$f(x) = x^2$ & 4 & 1 & 0 & 1 & 4 \\
\hline
\end{tabular}
\end{center}

\smallskip
The inputs \ans{$-2$} and \ans{2} both produce the output 4. Reversing
would send 4 to two places at once, so $f$ is \textbf{not} \ans{one-to-one} on all
real numbers --- and therefore not \ans{invertible} there.

\medskip
\textbf{The repair: restrict the domain.} Keep only $x \ge 0$. On that piece
the outputs never repeat, and

\begin{center}
$f^{-1}(x) = $ \ans{$\sqrt{x}$} \qquad with domain \ans{$x \ge 0$}
\end{center}

\medskip
\textbf{The graph of an inverse.} Every point $(a, b)$ on $y = f(x)$ becomes
the point $(b, a)$ on $y = f^{-1}(x)$ --- so the two graphs are mirror images
across the dashed line below.

\smallskip
\begin{center}
\begin{tikzpicture}[x=0.44cm, y=0.44cm]
  \draw[linegray, very thin, step=1] (-5,-5) grid (5,5);
  \draw[->, thick] (-5.4,0) -- (5.7,0) node[below right] {\small $x$};
  \draw[->, thick] (0,-5.4) -- (0,5.7) node[above] {\small $y$};
  \foreach \x in {-4,-2,2,4} \node[below, font=\tiny] at (\x,0) {\x};
  \foreach \y in {-4,-2,2,4} \node[left, font=\tiny] at (0,\y) {\y};
  \draw[dashed, thick, slate] (-5,-5) -- (5,5);
  \node[font=\tiny, slate, above left] at (4.6,4.6) {$y = x$};
  \draw[very thick, plum] (-3,-5) -- (2,5);
  \draw[very thick, goldacc] (-5,-3) -- (5,2);
  \fill[plum] (0,1) circle (2.2pt);
  \fill[goldacc] (1,0) circle (2.2pt);
  \fill[plum] (2,5) circle (2.2pt);
  \fill[goldacc] (5,2) circle (2.2pt);
\end{tikzpicture}
\end{center}
{\small\textcolor{plum}{\textbf{plum}} $= f(x) = 2x + 1$ \qquad
 \textcolor{goldacc}{\textbf{gold}} $= f^{-1}(x) = \dfrac{x-1}{2}$ \qquad
 \textcolor{slate}{\textbf{dashed}} $= y = x$}

\medskip
\begin{itemize}[itemsep=4pt]
  \item The point $(0, 1)$ on $f$ corresponds to the point $($ \ans{1} ,
        \ans{0} $)$ on $f^{-1}$.
  \item The point $(2, 5)$ on $f$ corresponds to the point $($ \ans{5} ,
        \ans{2} $)$ on $f^{-1}$.
  \item The two graphs are reflections of each other over the line
        \ans{$y = x$}.
\end{itemize}
\end{notesbox}

\vspace{0.1in}

\boxguard
\begin{practicebox}
\begin{enumerate}[itemsep=8pt]
  \item Let $f(x) = 3x + 6$. Find $f^{-1}(x)$, then verify by computing
        $f^{-1}(f(x))$.

        \begin{work}
                      y &= 3x + 6 \\
                      x &= 3y + 6 \\
                  x - 6 &= 3y \\
          \frac{x-6}{3} &= y \\
          f^{-1}\big(f(x)\big) &= \frac{(3x + 6) - 6}{3} \\
                               &= \frac{3x}{3} \\
                               &= x
        \end{work}

        $f^{-1}(x) = $ \ans{$(x-6)/3$} \qquad
        $f^{-1}\big(f(x)\big) = $ \ans{$x$}

  \item A function $f$ satisfies $f(-4) = 7$.

        \begin{enumerate}[label=(\alph*), itemsep=4pt]
          \item $f^{-1}(7) = $ \ans{$-4$}
          \item In one sentence, why is $f^{-1}(7)$ not equal to
                $\dfrac{1}{7}$?

                \ansline{The $-1$ is a name, not an exponent --- no reciprocal is involved.}
        \end{enumerate}
\end{enumerate}
\end{practicebox}

\end{document}
